% Part III of the guide "Lean through Artform Calculus": the textbook's
% theorems in Lean. Listings are cut from ../RepositoryTour.lean and from the
% library under ../../textbook/ArtformTextbook by ../tools/extract_listings.py, and the
% book's statements are quoted from the inventory of the edition the
% repository follows. Written 2026-10-11 (Phase 5.1 of the repository's plan).
\chapter*{Part III. The textbook's theorems in Lean}
\addcontentsline{toc}{chapter}{Part III. The textbook's theorems in Lean}

The repository that holds this guide also writes the numbered statements of
the book in Lean. Most have proofs that Lean checks; the status table records
what is complete. The folder \lc{type-theory} holds a paper on the
role-based type theory the book uses, \emph{Five Roles for Artforms}, with a
Lean library that checks the paper's first part. Part~II wrote the
constructions of Chapter~1 in core Lean, in one file. The repository's
library is written with Mathlib, Lean's library of mathematics, and follows
the whole book. This part says what the repository proves and how to build
it, introduces the parts of Mathlib that the library uses, and shows how a
statement of the book is written in Lean. It then reads four statements
closely, from the book's words to their proofs, and ends with the checks the
repository makes on itself.

\section*{What the repository proves}
\addcontentsline{toc}{section}{What the repository proves}

The repository gives each of the book's \statustotal{} numbered statements
one of five states:
\begin{itemize}[leftmargin=1.4em]
\item \emph{planned}: no Lean statement has been written yet;
\item \emph{stated}: a Lean statement has been written, and it waits for review
  against the book, for a revision after review, or for a proof;
\item \emph{proved}: an independent review has approved the Lean
  statement, and Lean has checked a proof of it;
\item \emph{partial}: as proved, for a narrower statement, with the reason
  recorded;
\item \emph{informal}: not to be written in Lean, by decision, with the
  reason recorded.
\end{itemize}

A statement is proved in this sense when three things hold. The library has
a Lean proposition tied to the statement's label in the book. An independent
AI reviewer has
read that proposition beside the book and approved it as saying what the
book says. And Lean has checked a proof of it that uses no axiom beyond
the three that the last section of this part describes.

Lean checks the third point, and only the third: that the proof proves the
Lean proposition. It cannot check the second, which compares a text in
mathematical English with a text in Lean. Three things help the reviewer, and
any reader, to judge it. Every Lean statement has a \emph{fidelity note},
which lists each way it reads the book: a hypothesis made explicit, a case
generalized, a choice between two readings. Non-vacuity examples exhibit
inputs on which the hypotheses hold. And an approval records hashes of the
Lean text and exact terms, the inventoried book text, and the fidelity notes.
If these change later, the statement goes back for review.

When this guide was built, \statuscount{proved} statements were proved,
\statuscount{partial} were partial, \statuscount{stated} were stated and
\statuscount{planned} were planned.
The appendix lists all of them. The file \lc{textbook/STATUS.md} of the
repository is always current.

The review of 11 October 2026 returned six translations for revision:
3.32, 5.44, 6.3, 6.11, 13.8 and 13.9. Their proofs check, but their Lean
forms do not yet cover the book's full claims. Statements 12.2 and 13.10
have approved translations and still need proofs; 12.15 and 13.7 still need
translations. The two partial approvals, 3.43 and 7.4, cover only the scope
recorded in their fidelity notes. The review and its reasons are in
\lc{textbook/review/audit-2026-10-11.md}.

\section*{Getting and building it}
\addcontentsline{toc}{section}{Getting and building it}

The repository is on GitHub. To build it, install VS~Code and its Lean~4
extension, as in Part~I. The extension also installs \lc{elan}, which
installs the version of Lean that a project asks for, and \lc{lake},
Lean's build tool. Then, in a terminal:
\begin{shellcommands}
git clone https://github.com/skrivov/artform-calculus-textbook
cd artform-calculus-textbook
lake exe cache get
lake build
\end{shellcommands}
The first two lines copy the repository and enter it. The file
\lc{lean-toolchain} names the version of Lean, \leanversion, and \lc{elan}
installs it the first time \lc{lake} runs. \lc{lake exe cache get} fetches
Mathlib and downloads its compiled files, about 450\,MB, instead of
compiling Mathlib, which would take hours. Unpacked, they take about 8\,GB.
\lc{lake build} then compiles the library itself in a few minutes. It
prints a line for each file, and an error for a file that Lean refuses.

In VS~Code, open the folder of the repository. The extension checks each
file as it is opened and shows Lean's answers in the infoview, as in
Part~I. The Lean playground in the browser has Mathlib but not this
library, so the files of the repository cannot be tried there. The files
of Parts~I and~II can.

\section*{From core Lean to Mathlib}
\addcontentsline{toc}{section}{From core Lean to Mathlib}

Parts~I and~II used Lean's core language alone. The library uses Mathlib, a
library of mathematics written in Lean and kept by its community
\cite{Mathlib2020}. It has the integers modulo~$n$, finite sums, groups,
linear maps, measures and categories, each with its theorems. A file asks
for a part of Mathlib with \li!import!. The listings of this section are cut
from the file \lc{RepositoryTour.lean}, which imports the whole library. As
in Parts~I and~II, Lean checks every answer they show.

\paragraph{Finite types.}
Part~II gave each finite type a list of its values, the class \li!Listed!,
so that \li!decide! could run through them. Mathlib's class for the same
purpose is \li!Fintype!, and \li!Fintype.card! counts the values. The
library gives the cells of every rod an instance of it, so a glued rod can
be counted:
\leanlisting{tfintype}
The last line is the check of Part~I on squares and remainders. \li!decide!
works as it did there, with Mathlib's instances.

\paragraph{Finite sums.}
A \li!Finset! is a finite set of values of a type, and
\li!∑ k ∈ s, f k! is the sum of \li!f! over the set \li!s!.
\li!Finset.range 5! is $\{0,1,2,3,4\}$, so the first sum counts the ten
points of the tetractys again. The last example proves Gauss's formula of
Part~I for Mathlib's sums. \li!Finset.sum_range_succ! is a theorem of
Mathlib that splits off the last term of a sum over a range, and \li!ring!
proves the polynomial identity that is left.
\leanlisting{tsums}

\paragraph{The integers modulo $n$.}
\li!ZMod n! is the integers modulo $n$, with their addition and
multiplication. The toy electrodynamics of Chapter~6 counts modulo~$8$:
there $5+4$ is $1$, and $3-5$ is $6$. For positive $n$, \li!ZMod n! is finite, so \li!decide!
checks a statement for all its values. Eight equal steps bring every value
back, and a number modulo~$3$, the values of the square's law
(\cref{fig:law-square}), is $0$, $1$ or~$2$:
\leanlisting{tzmod}

\paragraph{Permutations.}
A bijection of a type to itself is a permutation, \li!Equiv.Perm!.
Permutations form a group, written as a product, and \li!f * g! is
$f\circ g$: first $g$, then $f$, as the last line checks by computation.
\li!finRotate 4! moves each of the numbers $0,1,2$ on by one and $3$ to~$0$,
as the tick of the square moves each value one corner on. Its fourth power
is the identity and its square is not:
\leanlisting{tperm}

\paragraph{Additive maps.}
A tally of Chapter~5 is additive: $T(p+q)=T(p)+T(q)$. Mathlib's type of
additive maps from \li!A! to \li!B! is \li!A →+ B!: a function together with
the proof that it is additive. A theorem about all additive maps, such as
\li!map_add!, then applies to every tally. The library writes the law of the
square and its total~$T$ this way. The tick moves each value one corner on,
and the total of $(1,2,0,1)$ is $4$, which is $1$ modulo~$3$:
\leanlisting{ttally}

\paragraph{Finding a theorem.}
Mathlib has about 180\,000 theorems, and the first question
in a proof is often whether the step needed is one of them. \li!exact?!
searches the theorems in scope for one that closes the goal, and prints the
call it found:
\leanlisting{texact}
\leanreply{texact}
Loogle (\url{https://loogle.lean-lang.org}) searches Mathlib by name and by
the shape of a statement. The documentation of Mathlib
(\url{https://leanprover-community.github.io/mathlib4_docs}) lists every
definition and theorem, with its source.

\section*{How the book is written in Lean}
\addcontentsline{toc}{section}{How the book is written in Lean}

The library has three kinds of declaration:
\begin{itemize}[leftmargin=1.4em]
\item A definition of the book is a definition of Lean in the namespace
  \li!Artform!: a structure, an inductive type or a function. The rods of
  Chapter~1 are one, the canvas of \cref{def:canvas} another.
\item A statement of the book is a definition of a proposition,
  \li!def … : Prop!, in the namespace \li!Artform.Statement!. It says what
  the book says, and proves nothing.
\item Its proof is a theorem whose type is exactly that proposition.
\end{itemize}
The attribute \li!@[book "label"]! ties a declaration to the statement or
definition of the book that carries that label. Labels are used, not
numbers, because the numbers change between editions of the book and the
labels do not. The command \li!#book! lists the declarations tied to a
label:
\leanlisting{tbook}
\leanreply{tbook}
The statement of Noether's theorem is
\lc{Artform.\allowbreak Statement.\allowbreak noether}, a proposition, and
its proof is the theorem \lc{Artform.\allowbreak noether}, whose type is
that proposition. Keeping the two apart lets the reviewer judge the
statement without reading the proof, and lets a proof be replaced without
touching what it proves.

Each chapter of the book has a folder of the library, \lc{Ch01} to
\lc{Ch13}. Its \lc{Defs} folder holds the chapter's definitions,
\lc{Statements.lean} or the \lc{Statements} folder its statements, and the
other files their proofs. The file \lc{textbook/manifest.toml} gives every
statement and definition of the book its Lean names, its status and its
fidelity note. \lc{textbook/STATUS.md} and the appendix of this guide are
made from it.

\section*{Four statements read closely}
\addcontentsline{toc}{section}{Four statements read closely}

The four statements below come from Chapters~1 and~5 of the book. Each is
shown as the book prints it and as the library states it, with a table that
reads one into the other, then with its proof or the outline of its proof,
and with its fidelity note. In the library's listings the documentation
comments are left out; the files in the repository have them.

\subsection*{The laws of a rod}

Part~II built rods by three rules, substitution being the third, and ordered
their cells by position. The library follows \cref{sec:rules-rods}, which
builds rods by two rules and defines their order by cases:
\leanlisting{lrods}
\li!cells! is the elimination rule of \cref{sec:rules-rods} written as a
definition by cases. A cell of $R_n$ is a number below $n$, and a cell of
$\glue(P,Q)$ is a tag with a cell of the rod that the tag names.
\li!before! says, case by case, when one cell comes before another: in
$R_n$ by the order of the numbers, and in $\glue(P,Q)$ every cell tagged $L$
before every cell tagged $R$, and two cells with one tag as in $P$ or
in~$Q$. The attribute \li!@[reducible]! lets Lean see through \li!cells!,
so that Mathlib's instances for \li!Fin n! serve the cells of~$R_n$.

\bookstatement{prop:rod-laws}
\leanlisting{lrodlaws}

\begin{center}
\small
\renewcommand{\arraystretch}{1.15}
\begin{tabular}{@{}>{\raggedright\arraybackslash}p{.42\textwidth} >{\raggedright\arraybackslash}p{.54\textwidth}@{}}
\toprule
Book & Lean\\ \midrule
$\cells(X)$ & \li!X.cells!\\
$a$ comes before $b$ & \li!X.before a b!\\
no cell comes before itself & \li!¬ X.before a a!\\
of two different cells one comes before the other & \li!X.before a b ∨ X.before b a!, when \li!a ≠ b!\\
\bottomrule
\end{tabular}
\end{center}

The proof follows the book's, by induction over the two rules:
\leanlisting{lrodlawsproof}
\li!induction X with! splits the goal into the two rules. For a glued rod it
supplies the three laws for $P$ and for $Q$, which \li!obtain! takes apart
into \li!irrP!, \li!transP!, \li!totP! and the same for~$Q$.
\li!refine ⟨?_, ?_, ?_⟩! leaves one goal for each law. Each
\li!rintro ⟨_ | _, a⟩! introduces a cell and splits it by its tag, so that
each law becomes one goal for each combination of tags: two for the first
law, eight for the second and four for the third. Cells with one tag use the
laws of $P$ or of $Q$. A cell tagged $L$ comes before a cell tagged $R$ by
\li!trivial!, and the remaining cases have \li!False! as a hypothesis, which
\li!elim! turns into any goal.

\emph{Fidelity note.} Rods are built by the two rules of
\cref{sec:rules-rods}, and their cells and order by its computation rules.

\subsection*{Reversing a tick of two steps}

\bookstatement{prop:two-step-reversal}
The states are any type \li!S!, and the two steps and the involution are
permutations of it, so that $S_0^{-1}$ and $S_1^{-1}$ exist. In the group of
permutations \li!f * g! is $f\circ g$:
\leanlisting{lreversal}

\begin{center}
\small
\renewcommand{\arraystretch}{1.15}
\begin{tabular}{@{}>{\raggedright\arraybackslash}p{.44\textwidth} >{\raggedright\arraybackslash}p{.52\textwidth}@{}}
\toprule
Book & Lean\\ \midrule
$F=S_1\circ S_0$ & \li!S₁ * S₀!\\
$R$ is an involution & \li!R * R = 1!\\
$R\circ S_0\circ R=S_0^{-1}$ & \li!R * S₀ * R = S₀⁻¹!\\
$I_0=R\circ S_0$, $I_1=S_1\circ R$ & \li!R * S₀!, \li!S₁ * R!\\
$I_0$ is an involution & \li!(R * S₀) * (R * S₀) = 1!\\
$I_0\circ F\circ I_0=F^{-1}=I_1\circ F\circ I_1$ & the last two lines\\
\bottomrule
\end{tabular}
\end{center}

The proof is the book's chain of rewrites:
\leanlisting{lreversalproof}
\li!intro! names the three hypotheses \li!hR!, \li!h₀! and \li!h₁!. The proof
first derives $S_0\circ R\circ S_0=R$ from \li!h₀!, and the same for $S_1$
(\li!k₀! and \li!k₁!). It uses two theorems of Mathlib about groups: an
element whose product with $x$ is $1$ is the inverse of $x$, and an
involution is its own inverse. \li!hF! then computes
$F^{-1}=R\circ S_0\circ S_1\circ R$. Each of the four claims is a
\li!calc!, a chain of equations, each with its reason: \li!rw [h₀]! rewrites
with an equation, and \li!simp only [mul_assoc]! regroups a product, which
the book does without saying so. The proof file ends with an example in
which the hypotheses hold and $R$ is not the identity: the two shears of
\cref{sec:reversible-shears} on $\mathbb Z/8\times\mathbb Z/8$, with the flip
$R(p,q)=(-p,q)$.

\emph{Fidelity note.} The two steps and the involution are permutations of
an arbitrary type of states.

\subsection*{Noether's theorem, and the square}

An affine law (\cref{def:affine-law}) acts on the tables of a set of cells
\li!C!, with values in \li!ZMod n!:
\leanlisting{laffine}
\li!M! is an additive bijection, \li!≃+! in Mathlib, and \li!b! is the
source. \li!toDetLaw! reads the affine law as a deterministic law
(\cref{def:tick}), whose symmetries and conserved readings are those of
\cref{def:law-symmetry}:
\leanlisting{ldetlaw}
Tallies, pairings and keeping a pairing are written with additive maps:
\leanlisting{ltally}
A pairing sends a table to a tally, so that \li!P w! is the tally
$T_w=\langle w,\cdot\rangle$ and \li!P w p! is $\langle w,p\rangle$.

\bookstatement{thm:noether}
\leanlisting{lnoether}

\begin{center}
\small
\renewcommand{\arraystretch}{1.15}
\begin{tabular}{@{}>{\raggedright\arraybackslash}p{.36\textwidth} >{\raggedright\arraybackslash}p{.60\textwidth}@{}}
\toprule
Book & Lean\\ \midrule
$F(p)$ & \li!L.tick p!\\
$F$ keeps the pairing & \li!L.Keeps P!\\
adding $w$ to every state is a symmetry & \li!L.toDetLaw.IsSymmetry (· + w)!\\
$T_w(F(p))=T_w(p)+\langle w,b\rangle$ & \li!P w (L.tick p) = P w p + P w L.b!\\
$T_w$ is conserved & \li!L.toDetLaw.Conserved (P w)!\\
$T_w$ for exactly one table $w$ & \li|∃! w, P w = T|\\
\bottomrule
\end{tabular}
\end{center}

\li!(· + w)! is the function $p\mapsto p+w$. The proof follows the book's.
For part~1, $\langle w,M(p)\rangle=\langle M(w),M(p)\rangle=\langle
w,p\rangle$. For part~2, a conserved tally has $T(b)=T(F(0))=T(0)=0$, so $M$
keeps it; the pairing is perfect, so $T=T_w$ for one $w$, and $M$ fixes that
$w$.

The law of the square moves each value one corner on, modulo~$3$:
\leanlisting{lsquare}
\li!squareShift! is the additive bijection $p\mapsto(i\mapsto p(i-1))$,
with its inverse and the proofs that the two undo each other.
\li!squareLaw! has source $0$, and \li!squareTotal! is the total of the four
values, with the proof that it is additive.

\bookstatement{cor:noether-square}
\leanlisting{lnoethersquare}
The first line says that the fixed tables are exactly the constant ones, and
\li!2 • squareTotal! is $2T$. The proof applies Theorem~5.6 to the square:
the tick keeps the pairing $\sum_ip(v_i)q(v_i)$, that pairing is perfect,
and a fixed table is constant. Its last step uses that a number modulo~$3$
is $0$, $1$ or $2$, which \li!decide! checks, as in the listing on \li!ZMod!
above.

\emph{Fidelity notes.} Theorem~5.6 is stated for every $n$. The case $n=0$,
in which \li!ZMod 0! is the integers, is included, so the Lean statement is
slightly more general than the book's. In Corollary~5.7 the corners are
\li!Fin 4!, numbered anticlockwise, and ``the three constant ones'' is
stated as: the fixed tables are exactly the constant tables.

\subsection*{Discrete Stokes on a canvas}

Chapters~5 and~6 use canvases of dimension two with a stored orientation.
The library writes \cref{def:canvas} for that case:
\leanlisting{lcanvas}
A canvas has finite types of vertices, edges and faces, whose values can be
compared. The six fields in square brackets say so, and Lean finds them by
itself. Each edge has a source and a target. Each face has its boundary
cycle, with the proof that it is closed. The cycle is a list of edges, each
marked \li!true! when the cycle runs along the edge's stored direction and
\li!false! when it runs against it. How a face reads an edge, and when the
faces are read anticlockwise, are defined from the cycles:
\leanlisting{linc}
\li!inc e f! counts the times the boundary of \li!f! runs along \li!e!, less
the times it runs against it: an entry of the book's incidence matrix. A
canvas is anticlockwise when no face runs along an edge twice, and two
different faces that share an edge read it in opposite directions, as faces
drawn in the plane and read anticlockwise do. The operators of
\cref{sec:difference-operators} follow, for a canvas \li!K! and values in a
commutative ring \li!R!:
\leanlisting{lcalculus}

\bookstatement{prop:discrete-stokes}
\leanlisting{lstokes}

\begin{center}
\small
\renewcommand{\arraystretch}{1.15}
\begin{tabular}{@{}>{\raggedright\arraybackslash}p{.48\textwidth} >{\raggedright\arraybackslash}p{.48\textwidth}@{}}
\toprule
Book & Lean\\ \midrule
$\curl D=0$, $\divg\curl^*=0$ & \li!K.curl (K.diff φ) = 0!, \li!K.div (K.curlStar B) = 0!\\
each face read anticlockwise & \li!K.Anticlockwise →!\\
the edges on exactly one face of $S$ & \li!(S.filter …).card = 1!\\
each read as that face reads it & \li!∑ f ∈ S, K.inc e f!\\
the edges with exactly one end in $U$, read leaving $U$ & source in \li!U! and target outside, less the reverse\\
\bottomrule
\end{tabular}
\end{center}

The proof of each part is a lemma of \lc{Ch05/Stokes.lean}, and the theorem
collects them:
\leanlisting{lstokesproof}
$\curl D=0$ telescopes: round a closed cycle, the differences
$\varphi(\text{end})-\varphi(\text{start})$ of its steps add up to zero.
Summation by parts exchanges two finite sums, and $\divg\curl^*=0$ follows
from it. Stokes exchanges the sums too, and then uses that the canvas is
anticlockwise: an edge on two faces of $S$ is read by them in opposite
directions, so its two terms cancel. The two squares of
\cref{sec:two-squares} are an anticlockwise canvas, which \li!decide!
checks:
\leanlisting{ltwosquares}
\lc{![0, 1, 2, 3, 1, 4, 5]} is the function on \li!Fin 7! with those values.
The vertices $v_0,\ldots,v_5$ are \li!0! to \li!5!, the edges
$e_1,\ldots,e_7$ are \li!0! to \li!6!, and the right face runs against the
shared edge $e_2$.

\emph{Fidelity note.} The book asks, for the whole proposition, that the
faces be read anticlockwise. Only Stokes uses it, and the Lean statement
assumes it there alone, which makes the other three parts stronger. For an
edge on exactly one face of $S$, the sum $\sum_{f\in S}\operatorname{inc}(e,f)$
has one term that is not zero, and it is how that face reads the edge.

\section*{How the repository checks itself}
\addcontentsline{toc}{section}{How the repository checks itself}

\li!#print axioms! lists the axioms a theorem depends on, through every
theorem it uses:
\leanlisting{taxioms}
\leanreply{taxioms}
These three are the axioms of Mathlib's mathematics. \li!propext! says that
two propositions that imply each other are equal. \li!Quot.sound! says that
a quotient identifies the values that its relation relates; the principle
\li!funext! of Part~II follows from it. \li!Classical.choice! picks a value
of any type that has one, which gives proof by contradiction. Every theorem
of the library uses at most these three, and a check run on every change
lists the axioms of every theorem and fails on any other.

Lean has a faster way to decide a statement, \li!native_decide!, which runs
compiled code instead of the kernel's own computation. Its result is trusted
only as far as the compiler is, and Lean records that trust as an extra
axiom of the theorem, so the check rejects it. The large computations of the
library use \li!decide +kernel! instead, which the kernel itself checks, on
data that a script computes beforehand. The sixty rotations of the
icosahedron (\cref{prop:mm-icosahedron}) are checked so.

An approval records three review hashes. The first is of the text the reviewer read:
the statement, and the definitions it uses, as Lean prints them. The second
is of their exact terms, as Lean stores them. The third records the book's
inventoried text and the fidelity notes, including whether the approval is
partial. A change to any of these sends the statement back for review.
The book's inventory does not record dependencies between its definitions
and statements, so a change to its text conservatively reopens every review.

GitHub runs the same commands on every change:
\begin{itemize}[leftmargin=1.4em]
\item \lc{lake build --wfail} compiles the library and the files of this
  guide, and fails on any warning; the axioms of every theorem are then
  checked;
\item \lc{scripts/Report.lean} lists every declaration tied to the book, and
  \lc{scripts/status.py --check} checks that each proof proves exactly its
  statement and that the status table is current;
\item \lc{scripts/check\_tree.py} checks that no library source uses \li!sorry!
  or another forbidden construct;
\item the guide's own check fails if this PDF is not built from the current
  sources, including the inputs to its generated status appendix.
\end{itemize}
Each release also runs \lc{leanchecker}, a second checker, separate from the
one that compiled the library, which reads every declaration again from the
compiled files.

\section*{Following a proof, and reporting a problem}
\addcontentsline{toc}{section}{Following a proof, and reporting a problem}

To follow a proof, open its file in VS~Code and put the cursor after a
tactic. The infoview shows the goal at that point, with its hypotheses.
Holding the control key (on a Mac, the command key) and clicking a name opens
its definition, in the library or in Mathlib, and hovering over a name shows
its type and documentation.

A Lean statement may seem to say less or more than the book, or the book may
seem wrong. The file \lc{textbook/FINDINGS.md} records what has been found
so far, and an issue on the repository's GitHub page reports something new.
A statement is never weakened to make a proof go through. One that cannot be
proved as the book prints it stays unproved, with the reason recorded, until
the book is corrected.
